\section{Beyond thermodynamics}\label{to:sec:beyond}

The same mechanism appears wherever a released system must stay correlated with what the device keeps. Its sharpest form is a statement about quantum mechanics alone, best seen on one task.

\paragraph{The delayed-choice task.} The task has two outputs, a qubit $B$ and a classical bit $a$. The observer supplies a question $x$, which chooses one of two bases, and the device must return a uniformly random answer $a$ such that, given $x$ and $a$, the qubit $B$ is the corresponding basis state; for the BB84 task the bases are the computational and Hadamard bases. In batch the question arrives before $B$ is released. In real time it arrives afterwards, and the device never touches $B$ again. A Bell pair realizes the real-time task: release one half, and measure the retained half in the requested basis.

\begin{lemma}[Release lemma]\label{to:lem:release}
Suppose a qubit $B$ is released, and afterwards a memory $M$, which holds everything the device keeps at the release, must produce an answer $a$ such that, for either of two questions $x$ chosen after the release, the conditional states of $B$ are pure in a basis that depends on $x$. If the two bases do not commute, every exact realization has $-S(B|M)=1$. If each mismatch probability is at most $\varepsilon$, in particular at strategy error $\varepsilon$, then $-S(B|M)\ge1-g\bigl(\sqrt\varepsilon\,(1+2/\sin2\theta)\bigr)$ with $g(d)=2d+(1+d)\,h_2\bigl(d/(1+d)\bigr)$, where $\theta$ is the angle between the bases, independently of the dimension of $M$.
\end{lemma}

The exact statement is closely related to known results. The assemblages of the two questions are extremal for every nonzero angle, and Cope showed that every state realizing an extremal assemblage has entanglement of formation equal to the entropy of its reduced state, here one ebit~\cite[Thm.~1]{Cope21}; for the BB84 pair it also follows from one-sided self-testing of the ebit~\cite{SupicHoban16}, and the uncertainty relation with quantum memory gives $-S(B|M)\ge1-2h_2(\varepsilon)$~\cite{BCCRR10}. What the lemma adds is the form, a statement about whatever memory outlives a release, usable in every resource theory, and a robust bound that does not depend on the dimension of $M$.

Each resource theory converts $-S(B|M)$ into its own currency.
\begin{itemize}[leftmargin=*]
\item \emph{Entanglement under one-way communication from the party that answers to the party that releases.} For the delayed-choice BB84 task, holding the qubit costs nothing, while releasing it first costs exactly one ebit per use, extensively. The uncertainty relation with quantum memory~\cite{BCCRR10}, with Fano's inequality, gives $-S(B|M)\ge1-2h_2(\varepsilon)$; squashed entanglement~\cite{ChristandlWinter04} is at least $-S(B|M)$, by strong subadditivity; and its monogamy~\cite[Thm.~8]{KoashiWinter04} charges every release, so $n$ uses need $n[1-2h_2(\varepsilon)]$ ebits. By the release lemma the rate is the same for every tilt $\theta\in(0,\pi/2)$. Two-way communication removes the cost, because the releasing party can keep the partner of $B$ and report the answer.
\item \emph{Coherence under incoherent operations}~\cite{BCP14,WinterYang16}. For the delayed choice between the computational basis and a basis tilted by $\theta$, real time costs one coherent bit per use and batch costs $h_2(\sin^2\theta)$, the coherence cost~\cite{WinterYang16} of the tilted basis states, so the surcharge is $1-h_2(\sin^2\theta)$, $0.399$ at $\theta=\pi/8$. The surcharge approaches one bit as $\theta\to0$: for every fixed $\theta>0$ the real-time cost at vanishing error stays one coherent bit, while the batch cost $h_2(\sin^2\theta)$ tends to zero. It vanishes for mutually unbiased bases. A two-step process built from a CNOT and a dephasing of the first output in the $X$ basis, which is a free mixture of $I$ and $X$, costs nothing in batch and one coherent bit per episode in real time; $N$ episodes at error $\varepsilon$ consume at least $N[1-h_2(\varepsilon)]$ coherent bits net (Appendix~\ref{to:app:numbers}).
\end{itemize}
Blank ancillas in a computational-basis state, discarding, and reset to a computational-basis state are free in both theories. In real time the answering party's messages reach the releasing party only after $B$ has left, so in these examples the real-time entanglement cost is the entanglement cost of the task's steering assemblage, which Cope defined~\cite{Cope21}. In the resource theory of steering, communication may run only from the party that holds the steered system, since communication towards it could create steering~\cite{GallegoAolita15}. Here it runs towards that party: in batch this is what makes the task free, and in real time the release order is what keeps it from creating steering. What is new is this comparison with batch, the surcharge for releasing rather than holding the same process. A per-use timing cost does not need forgetting to be priced; it needs only that the cheapest global implementation cannot be split, at the same cost, into stages that respect the release order.

\paragraph{Towards a general law.} In the examples above, the gross cost of a schedule is the smoothed regularized formation cost of one common source together with a free causal decoder, and the time-order surcharge is the difference between its real-time and batch values. We expect the same under general formation hypotheses, for instance for one-way communication whenever all released outputs leave before the first question; we do not state or prove them here. For a single release this extends the entanglement cost of steering assemblages~\cite{Cope21}. For channels, the entanglement cost of simulating uses called in sequence by adaptive testers, rather than in parallel, was studied by Wilde~\cite{Wilde18}. The smoothing matters: the entanglement of formation of an assemblage is not continuous, and Cope expects its plain regularization to differ from the cost~\cite{Cope21}. Cope and Uola bound the regularized smoothed Schmidt measure of an assemblage below by the regularized entanglement of formation of its smoothed copies and above by that of its copies themselves; identifying the bounds with the entanglement cost would need a continuity that fails for general assemblages, and they leave open whether the two bounds can differ~\cite[Sec.~VI.A, Eq.~(22)]{CopeUola22}. A general equality here meets the same gap. Removing the smoothing, crediting returned resources and factorizing interleaved releases remain open. Three quantities must be kept apart: consumption, a rate; working capital, the largest running deficit, which can be refilled by resources that arrive later; and sublinear irreversible overhead. The naive statement that the surcharge is the correlation forced across the release cut fails: in the tilted coherence task the forced correlation is one full bit, while the surcharge is $1-h_2(\sin^2\theta)$, because part of the correlation is later delivered to the observer as the output it asked for. That part must be credited.

\paragraph{No single currency.} The delayed-choice task certifies exactly one qubit of necessary quantum retention: a classical record across the cut, an entanglement-breaking memory in the sense of~\cite{RBL18}, cannot realize it, since $S(B|J)\ge0$ for classical $J$, and $-S(B|M)\le\log_2\dim M$ forces a qubit. This requirement is the same in every theory. Its price is not (Table~\ref{to:tab:currency}).

\begin{table}[tbp]
\centering
\small
\begin{tabular}{lccc}
\toprule
Theory & Real time & Batch & Order price\\
\midrule
Thermodynamics at zero Hamiltonian & $2\,\kB T\ln2$ & $\kB T\ln2$ & $\kB T\ln2$\\
Entanglement, one-way communication & $1$ ebit & $0$ & $1$ ebit\\
Coherence, incoherent operations & $1$ & $1$ & $0$\\
Magic, stabilizer operations & $0$ & $0$ & $0$\\
\bottomrule
\end{tabular}
\caption{The delayed-choice BB84 task in four resource theories. In entanglement theory communication is one-way, from the party that answers to the party that releases. Coherence is counted in coherent bits; under stabilizer operations a Bell pair and Pauli measurements are free, so the task costs nothing in either schedule.}\label{to:tab:currency}
\end{table}

In the thermodynamic row the output registers are taken from the reservoir, so no blank register is free, and the setting register is read-only and returned; with free blank registers both thermodynamic entries fall by $2\,\kB T\ln2$, so that batch then extracts $\kB T\ln2$, and the order price is unchanged (Appendix~\ref{to:app:numbers}). In thermodynamics the Bell pair that realizes the task in real time is worth two bits of work; one bit is delivered to the observer as the output, and one is the price of order. The order prices are not proportional across tasks. In the entanglement theory of Table~\ref{to:tab:currency}, with one-way communication from the party that answers to the party that releases, the BB84 task and the tilted task both carry an order price of one ebit; in coherence theory their order prices are $0$ and $0.399$. With rows for entanglement and coherence and columns for the BB84 and tilted tasks, the order prices form the matrix
\[
\begin{pmatrix}1&1\\0&0.399\end{pmatrix},
\]
which has rank two, so no amount of order per task multiplied by an exchange rate per theory reproduces it. Even an ordering of tasks by amount of order fails in these two theories: the BB84 task has order price $1$ in that entanglement theory and $0$ in coherence theory, while the CNOT process above, with its first slot at the party that answers and its second at the party that releases, has order price $0$ in that entanglement theory, since in either schedule the first party can measure its input in the $X$ basis, which is the required dephasing, and send the outcome to the second, who applies the corresponding phase to the second input; in coherence theory its order price is $1$. The placement is forced, not chosen: the process signals from its first slot to its second, so under one-way communication from the party that answers to the party that releases it can be realized at all only with its first slot at the party that answers; for this process the two names just mean the sending and the receiving party. What the delayed-choice examples share is the requirement, one retained Bell pair. When the answering measurement is free, as in every row of the table, real time pays that Bell pair's price in the theory, and the order price is that price minus what the task would cost in batch. The CNOT process needs no retained Bell pair: measuring its first input in the $X$ basis and keeping the outcome as a classical record realizes it in real time, so it is classical-retention in the sense of the next paragraph, and its coherence price is a price of coherent correlation, not of quantum retention. What replaces a currency is a budget region for each scheduled process, consumption, capital, reference and returns, which each theory values differently.

\paragraph{Order cannot be manufactured.} Call a process classical-retention if it can be realized with only classical records kept across each release cut. The property concerns what is retained across each cut, not whether every channel realized during a visit is entanglement-breaking: a carrier passed through unchanged and released before the cut leaves no quantum record across it, so on that schedule the identity channel is classical-retention, although it is not entanglement-breaking. For a memory link with its input before the cut and its output after it, classical retention gives exactly an entanglement-breaking channel, a free object of the resource theory of quantum memories~\cite{RBL18}; for many cuts the classical-retention processes are the processes with classical memory of the multi-time resource theories of Berk et al.~\cite{Berk21}. As in those theories, any contractive distance to the classical-retention processes, the strategy-norm distance included, cannot increase under schedule-respecting adapters that themselves keep only classical records across release cuts~\cite[Thm.~1]{Berk21}, so order cannot be manufactured by such adapters. Adapters that keep quantum memory are not free in this sense: an adapter that receives both halves of a Bell pair at once, releases one and holds the other across the release turns that source into the BB84 task. The surcharge itself is monotone only under conversions that preserve batch credits: a source that releases one half of a Bell pair early and delivers the other half late has surcharge $0$, and an adapter that keeps nothing, measuring the late half in the basis $x$, turns it into the BB84 task, whose surcharge is $1$.

% =====================================================================
